\documentclass[a4paper,11pt]{article}
\usepackage[T2A]{fontenc}
\usepackage[utf8]{inputenc}
\usepackage[russian]{babel}
\usepackage{amsmath,amssymb,mathtools}
\usepackage{helvet}
\renewcommand{\familydefault}{\sfdefault}
\usepackage{icomma}
\usepackage{geometry}
\geometry{left=22mm,right=22mm,top=20mm,bottom=21mm}
\usepackage{microtype}
\usepackage{tikz}
\usetikzlibrary{calc,arrows.meta,positioning,shapes.geometric}
\usepackage{booktabs,tabularx,array}
\usepackage{enumitem}
\usepackage{xcolor}
\usepackage{fancyhdr}
\usepackage{setspace}
\usepackage{hyperref}
\usepackage{needspace}
\definecolor{accent}{HTML}{008080}
\definecolor{darkaccent}{HTML}{004D4D}
\definecolor{textgray}{HTML}{333333}
\definecolor{softgray}{HTML}{777777}
\hypersetup{colorlinks=true,linkcolor=darkaccent,urlcolor=darkaccent}
\setstretch{1.26}
\setlength{\parindent}{0pt}
\setlength{\parskip}{5pt}
\setlist[itemize]{topsep=3pt,itemsep=3pt,leftmargin=6mm}
\setlist[enumerate]{topsep=4pt,itemsep=8pt,leftmargin=8mm}
\pagestyle{fancy}
\fancyhf{}
\fancyhead[L]{\footnotesize\color{softgray}Математика · Оптимизация}
\fancyhead[R]{\footnotesize\color{softgray}Тимофей Васильченко}
\fancyfoot[C]{\small\color{softgray}\thepage}
\renewcommand{\headrulewidth}{0.25pt}
\renewcommand{\footrulewidth}{0pt}
\setlength{\headheight}{14pt}
\newcommand{\topic}[1]{\Needspace{5\baselineskip}\section*{\color{darkaccent}#1}\vspace{-2mm}}
\newcommand{\key}[1]{\textcolor{darkaccent}{\textbf{#1}}}
\newcommand{\exercise}[2]{\item \textbf{#1}\quad #2}
\begin{document}
\thispagestyle{empty}
\begin{center}
\vspace*{23mm}
{\Large\color{darkaccent}МАТЕМАТИКА}\\[25mm]
{\Huge\color{textgray}Чек-лист}\\[13mm]
{\LARGE\color{accent}Задачи на оптимизацию}\\[3mm]
{\Large\color{accent}и производная сложной функции}\\[22mm]
{\large Подготовка к ЕГЭ по профильной математике}\\[43mm]
{\Large Тимофей Васильченко}\\[17mm]
{\large 14 ноября 2026 г.}\\[20mm]
{\normalsize Лёвин Артём Александрович}\\[2mm]
{\normalsize эксклюзивно для Тимофея Васильченко}
\end{center}
\vfill
\begin{center}
\begin{tikzpicture}[x=1cm,y=1cm]
\draw[darkaccent,line width=.8pt] (-5,0) .. controls (-3.3,.9) and (-1.8,-.65) .. (0,0) .. controls (1.9,.7) and (3.2,-.5) .. (5,0);
\fill[accent] (0,0) circle (1.7pt);
\end{tikzpicture}
\end{center}
\clearpage

\topic{1. Главная идея задач на оптимизацию}
\key{Оптимизация} — поиск наименьшего или наибольшего допустимого значения величины: расходов, прибыли, цены, налогов, времени.

\begin{itemize}[label=$\square$]
\item Понимаю, \textbf{какую величину нужно оптимизировать} и в каких единицах она измеряется.
\item Выбираю переменную и выписываю её \textbf{допустимые значения}.
\item Составляю \textbf{целевую функцию} $F(x)$, зависящую от одной переменной.
\item Нахожу $F'(x)$ и решаю уравнение $F'(x)=0$.
\item Проверяю знак производной, границы отрезка и ограничения.
\item Если переменная целочисленная, сравниваю значения функции в подходящих целых точках.
\item Возвращаюсь к тексту и записываю ответ с единицами измерения.
\end{itemize}

\key{Признак экстремума по знаку производной.}
Если $F'$ меняет знак с «$-$» на «$+$», получаем минимум; с «$+$» на «$-$» — максимум. Одного равенства $F'(x)=0$ недостаточно для окончательного вывода.

\begin{center}
\begin{tabular}{@{}lll@{}}
\toprule
Знак $F'$ слева & Знак $F'$ справа & Вывод\\
\midrule
$-$ & $+$ & локальный минимум\\
$+$ & $-$ & локальный максимум\\
Одинаковый & одинаковый & смены монотонности нет\\
\bottomrule
\end{tabular}
\end{center}

\key{Контроль границ.} На закрытом отрезке необходимо сравнить все внутренние критические точки и оба конца отрезка. В задаче с целым количеством объектов дополнительно учитывается дискретность допустимых значений.

\topic{2. Производная сложной функции}
Если $y=[u(x)]^n$, то
\[
\boxed{y'=n[u(x)]^{n-1}\cdot u'(x).}
\]
Сначала дифференцируем \textbf{внешнюю степень}, затем умножаем на производную \textbf{внутреннего выражения}.

\key{Пример с доски.} Пусть $f(x)=(x^2-x-6)^{10}$. Обозначим $u(x)=x^2-x-6$. Тогда
\[
f'(x)=10(x^2-x-6)^9(2x-1).
\]
\clearpage
Для определения стационарных точек разложим внутренний трёхчлен:
\[
x^2-x-6=(x+2)(x-3),\qquad
f'(x)=10(x+2)^9(x-3)^9(2x-1).
\]
Отсюда $x=-2$, $x=\tfrac12$, $x=3$. Знаки производной последовательно равны
\[
\begin{array}{c|cccc}
\text{Промежуток} & (-\infty,-2)&(-2,\tfrac12)&(\tfrac12,3)&(3,+\infty)\\\hline
f'(x)&-&+&-&+
\end{array}
\]
Значит, $x=-2$ и $x=3$ — точки минимума (значение $f=0$), а $x=\tfrac12$ — точка локального максимума.

\key{Особое внимание:} множитель $(2x-1)$ терять нельзя. Производная от $(u(x))^{10}$ равна $10u^9u'$, а не просто $10u^9$.

\topic{3. Распределение людей между объектами}
\key{Сюжет занятия.} Имеются 24 рабочих. Если на первом объекте работают $t$ человек, расходы на их оплату равны $4t^2$ у.е.; на втором работают $24-t$ человек, расходы равны $(24-t)^2$ у.е.

\textbf{Шаг 1. Модель и ограничения:}
\[
S(t)=4t^2+(24-t)^2=5t^2-48t+576,
\quad t\in\{0,1,\ldots,24\}.
\]
\textbf{Шаг 2. Непрерывный ориентир:}
\[
S'(t)=10t-48=0\quad\Longrightarrow\quad t=4{,}8.
\]
До $4{,}8$ функция убывает, после — возрастает. Количество работников должно быть целым.

\textbf{Шаг 3. Проверка соседних целых значений:}
\[
S(4)=4\cdot16+20^2=464,
\qquad S(5)=4\cdot25+19^2=461.
\]
\key{Ответ:} 5 человек на первом объекте и 19 — на втором. Минимальные расходы составят $461$ у.е.

\topic{4. Распределение выпуска между заводами}
\key{Сюжет занятия.} Для изготовления $t$ приборов на первом заводе требуется $4t^3$ рабочих часов, на втором для изготовления $20-t$ приборов — $(20-t)^3$ часов. Стоимость одного часа — 1 тыс. руб.

\textbf{Целевая функция:}
\[
S(t)=4t^3+(20-t)^3,
\qquad t\in\{0,1,\ldots,20\}.
\]
\clearpage
\textbf{Поиск минимума на вещественном промежутке:}
\begin{align*}
S'(t)&=12t^2-3(20-t)^2,\\
S'(t)=0 &\iff 4t^2=(20-t)^2.
\end{align*}
Так как $0\le t\le20$, обе величины $2t$ и $20-t$ неотрицательны:
\[
2t=20-t\quad\Longrightarrow\quad t=\frac{20}{3}\approx6{,}67.
\]
\textbf{Проверка целых вариантов:}
\[
S(6)=4\cdot6^3+14^3=3608,
\qquad S(7)=4\cdot7^3+13^3=3569.
\]
\key{Ответ:} 7 приборов на первом заводе, 13 — на втором; минимальная оплата — $3569$ тыс. руб.

\key{Общее правило двух задач.} Производная помогает локализовать оптимум, но в прикладных задачах количество людей и приборов должно быть целым. При нескольких действительных стационарных точках проверяют все допустимые кандидаты.

\topic{5. Минимальная цена при заданной прибыли}
\key{Сюжет занятия.} Производство $x$ тыс. единиц продукции стоит в год
\[
q(x)=0{,}5x^2+x+7\quad\text{млн руб.}
\]
Цена единицы продукции $p$ тыс. руб., выручка $px$ млн руб. За три года требуется получить не менее 75 млн руб. прибыли при неизменных ежегодных показателях.

\textbf{Условие прибыли:}
\[
3\bigl(px-q(x)\bigr)\ge75
\iff px-(0{,}5x^2+x+7)\ge25.
\]
При $x>0$ выразим необходимую цену:
\[
p\ge\frac{25+0{,}5x^2+x+7}{x}
=0{,}5x+1+\frac{32}{x}.
\]
Достаточно найти минимум функции $P(x)=0{,}5x+1+32/x$ при $x>0$:
\[
P'(x)=0{,}5-\frac{32}{x^2}=0
\quad\Longrightarrow\quad x^2=64\quad\Longrightarrow\quad x=8.
\]
Знак $P'$ меняется с «$-$» на «$+$», поэтому $P(8)=4+1+4=9$.

\key{Проверка:} выручка $8\cdot9=72$ млн руб. в год, издержки $0{,}5\cdot64+8+7=47$ млн руб.; за три года прибыль $3(72-47)=75$ млн руб.

\key{Ответ:} минимальная цена — $9$ тыс. руб. за единицу продукции.
\clearpage

\topic{6. Максимум налоговых поступлений}
\key{Сюжет занятия.} Спрос на товар задан формулой $Q=20000-P$ ($Q$ — число единиц, $P$ — цена в рублях). Издержки фирмы равны $6000Q+4\,000\,000$ руб. Государство устанавливает налог $t$ рублей на каждую единицу, $0\le t\le10000$.

\textbf{Шаг 1. Решение фирмы при фиксированном налоге.}
Из спроса $P=20000-Q$, следовательно прибыль фирмы:
\begin{align*}
K(Q,t)&=PQ-(6000Q+4\,000\,000)-tQ\\
&=-Q^2+(14000-t)Q-4\,000\,000.
\end{align*}
При заданном $t$ фирма выбирает $Q$ для максимума $K$:
\[
K'_Q(Q,t)=14000-t-2Q=0
\quad\Longrightarrow\quad Q(t)=7000-\frac t2.
\]
На интервале $0\le t\le10000$ получаем $2000\le Q(t)\le7000$, то есть ограничения спроса соблюдаются. Поскольку коэффициент при $Q^2$ отрицательный, это максимум прибыли фирмы.

\textbf{Шаг 2. Решение государства.}
Налоговые поступления зависят и от ставки, и от выпуска фирмы:
\[
N(t)=t\,Q(t)=7000t-\frac{t^2}{2},
\qquad N'(t)=7000-t.
\]
Максимум на $[0,10000]$ достигается при $t=7000$. Тогда
\[
Q=7000-\frac{7000}{2}=3500,
\qquad N=7000\cdot3500=24\,500\,000\ \text{руб.}
\]
\key{Ответ:} налог $7000$ руб. за единицу, поступления $24{,}5$ млн руб.

\key{Главная ловушка.} Если ставка налога увеличивается, фирма меняет объём производства. Поэтому выражение $tQ$ нужно оптимизировать с учётом зависимости $Q=Q(t)$.

\topic{7. Типичные ошибки и самопроверка}
\begin{itemize}[label=$\square$]
\item Различаю \textbf{выручку}, \textbf{затраты}, \textbf{прибыль} и \textbf{налоговые поступления}.
\item При делении не забываю условия $x\ne0$ и знак $x$, если делю неравенство.
\item Выделяю \textbf{внутреннюю функцию} при дифференцировании степени.
\item Проверяю \textbf{знак производной} по обе стороны критической точки.
\item Учитываю границы допустимого интервала, в том числе налоговые пределы.
\item После непрерывного поиска учитываю \textbf{целочисленность} работников и изделий.
\item Проверяю единицы: тыс. руб. $\times$ тыс. единиц $=$ млн руб.
\item В многоэтапной задаче сначала нахожу \textbf{реакцию фирмы}, затем оптимум государства.
\end{itemize}
\clearpage

\thispagestyle{fancy}
\begin{center}
{\Large\color{darkaccent}\bfseries Алгоритм: как решать задачу на оптимизацию}
\end{center}
\vspace{7mm}
\begin{center}
\begin{tikzpicture}[
 >=Stealth, node distance=10mm,
 start/.style={ellipse,draw=darkaccent,thick,minimum width=38mm,minimum height=9mm,align=center,fill=white},
 action/.style={rectangle,rounded corners=2mm,draw=accent,thick,text width=114mm,minimum height=12mm,align=center,inner sep=4pt},
 decision/.style={diamond,draw=darkaccent,thick,aspect=2.8,inner xsep=6mm,inner ysep=3mm,align=center,text width=37mm},
 branch/.style={rectangle,rounded corners=1.8mm,draw=accent,thick,text width=51mm,minimum height=14mm,align=center,inner sep=3pt},
 link/.style={->,thick,draw=darkaccent}
]
\node[start] (a) {Начало};
\node[action,below=of a] (b) {Назовите целевую величину и выберите переменную};
\node[action,below=of b] (c) {Установите область допустимых значений и единицы измерения};
\node[action,below=of c] (d) {Выразите целевую функцию $F(x)$ через одну переменную};
\node[action,below=of d] (e) {Найдите $F'(x)$, критические точки и характер монотонности};
\node[decision,below=12mm of e] (f) {Переменная целая?};
\node[branch,below left=17mm and 3mm of f] (g) {Да: проверьте соседние целые значения и границы};
\node[branch,below right=17mm and 3mm of f] (h) {Нет: проверьте критические точки и границы};
\node[action,below=28mm of f] (i) {Сравните допустимые значения $F$ и выберите минимум/максимум};
\node[start,below=of i] (j) {Запишите ответ с единицами};
\draw[link] (a)--(b);\draw[link] (b)--(c);\draw[link] (c)--(d);
\draw[link] (d)--(e);\draw[link] (e)--(f);
\draw[link] (f.south west) -| node[pos=.25,above left,fill=white,font=\small]{Да} (g.north);
\draw[link] (f.south east) -| node[pos=.25,above right,fill=white,font=\small]{Нет} (h.north);
\draw[link] (g.south) |- (i.west);
\draw[link] (h.south) |- (i.east);
\draw[link] (i)--(j);
\end{tikzpicture}
\end{center}
\clearpage

\topic{Тренировочный блок · Путь А}
\textcolor{softgray}{10 заданий, от базовых к сложным. Во всех задачах обоснуйте наличие минимума или максимума. Для прикладных задач укажите ограничения и единицы.}
\begin{enumerate}[label=\textbf{\arabic*.}]
\item Найдите наименьшее значение функции $F(x)=3x^2-18x+5$ на всей числовой прямой.

\item Найдите производную $f(x)=(x^2-3x-4)^5$ и все точки, в которых $f'(x)=0$.

\item Найдите наибольшее значение $R(x)=20x-x^2$ при $0\le x\le20$.

\item 18 рабочих распределяют между двумя объектами. Если на первом работают $t$ человек, расходы составляют $3t^2$ у.е., а на втором — $(18-t)^2$ у.е. Найдите все распределения с минимальными расходами. Число людей целое.

\item На первом заводе изготовление $t$ приборов требует $2t^3$ часов, на втором изготовление остальных $16-t$ приборов требует $(16-t)^3$ часов. Оплата часа одинаковая. Как распределить 16 приборов с минимальными затратами? Число приборов целое.

\item Найдите наименьшее значение $P(x)=\dfrac{x}{4}+1+\dfrac{16}{x}$ при $x>0$ и значение $x$, при котором оно достигается.
\clearpage
\item Выпуск $x$ тыс. изделий обходится в $C(x)=0{,}5x^2+2x+10$ млн руб. в год. Цена изделия $p$ тыс. руб. За два года при одинаковом выпуске и цене прибыль должна составить не менее 44 млн руб. Найдите минимальную допустимую цену и соответствующий объём выпуска; $x>0$.

\item Спрос задан как $Q=12000-P$, где $Q$ — количество единиц товара, $P$ — цена в рублях. Затраты фирмы — $2000Q+1\,000\,000$ руб. Налог $t$ руб. взимается с каждой единицы, $0\le t\le6000$. Фирма выбирает выпуск для максимума собственной прибыли. Найдите ставку $t$, обеспечивающую максимальные налоговые поступления, и величину этих поступлений.

\item Спрос задан как $Q=16000-P$. Издержки фирмы — $3000Q+1\,000\,000$ руб. Налог на единицу $t$ ограничен условием $0\le t\le6000$. Фирма максимизирует прибыль выбором $Q$. При какой ставке налоговые поступления наибольшие? Найдите их сумму. Обязательно проверьте границы для $t$.

\item Цена $p$ измеряется в тыс. руб., спрос $Q=100-2p$ — в тыс. единиц, $0\le p\le50$. Затраты на выпуск $Q$ тыс. единиц равны $C(Q)=\dfrac{Q^2}{4}+20Q+40$ млн руб. Найдите цену, при которой прибыль фирмы максимальна, оптимальный выпуск и максимальную прибыль.
\end{enumerate}
\clearpage

\topic{Ответы к тренировочному блоку}
\textcolor{softgray}{Для самопроверки приведены итоговые результаты. Сначала выполните задания самостоятельно.}
\renewcommand{\arraystretch}{1.43}
\begin{tabularx}{\linewidth}{@{}>{\bfseries}p{12mm}X@{}}
\toprule
№ & Ответ \\
\midrule
1 & $x=3$, $F_{\min}=-22$.\\
2 & $f'(x)=5(x^2-3x-4)^4(2x-3)$; $x=-1$, $\frac32$, $4$.\\
3 & $x=10$, $R_{\max}=100$.\\
4 & $t=4$ (4 и 14 рабочих) или $t=5$ (5 и 13 рабочих); минимум $244$ у.е.\\
5 & $t=7$ (7 и 9 приборов); минимум $1415$ часов.\\
6 & $x=8$, $P_{\min}=5$.\\
7 & $x=8$ тыс. изделий; $p_{\min}=10$ тыс. руб.\\
8 & $t=5000$ руб.; $Q=2500$; максимум $12\,500\,000$ руб.\\
9 & $t=6000$ руб. (правая граница); $Q=3500$; максимум $21\,000\,000$ руб.\\
10 & $p=40$ тыс. руб.; $Q=20$ тыс. единиц; прибыль $260$ млн руб.\\
\bottomrule
\end{tabularx}
\vspace{14mm}
\key{Итоговая памятка}
\begin{center}
\begin{tabularx}{\linewidth}{@{}p{0.34\linewidth}X@{}}
\toprule
Ситуация & Что делать \\
\midrule
Количество работников или изделий & Оптимизировать на непрерывном отрезке, затем сравнить допустимые целые варианты.\\
Минимальная цена при нужной прибыли & Выразить цену из условия прибыли и найти её минимум.\\
Налог на каждую единицу товара & Сначала найти выпуск фирмы $Q(t)$, затем максимизировать $tQ(t)$.\\
Степень от выражения & Применить правило $[u(x)^n]'=nu(x)^{n-1}u'(x)$.\\
\bottomrule
\end{tabularx}
\end{center}
\vfill
{\footnotesize\color{softgray}Лёвин Артём Александрович · эксклюзивно для Тимофея Васильченко}
\end{document}
